\documentclass{article}
\PassOptionsToPackage{numbers,sort&compress}{natbib}
\usepackage[preprint]{neurips_2026}
\usepackage[utf8]{inputenc}
\usepackage[T1]{fontenc}
\usepackage{amsmath}
\usepackage{microtype}
\usepackage[hidelinks]{hyperref}

% Author information.
\newcommand{\AuthorName}{Lin Zihan}
\newcommand{\StudentID}{50027371}
\title{Photo Editing with Controllable\\Styles, Filters, and Composition}
\author{Name: \AuthorName\quad Student ID: \StudentID}
\hypersetup{pdftitle={Photo Editing with Controllable Styles, Filters, and Composition},pdfsubject={DSAA 2012 Project Proposal}}

\begin{document}
\maketitle
\section{Project Goal}

I want to build a photo editor that turns an input image into a chosen style, filter, and composition while keeping its main subject recognizable. Example looks include warm film-like colors, muted vintage tones, and collage posters with off-center subjects. I will choose a few online reference looks and document their sources and visible popularity signals.

\section{How It Will Work}

Segmentation will separate the subject from the background. Cropping or repositioning will adjust subject placement and empty space; color filters will change the overall mood. Reference-guided styling \citep{wang2024instantstyle} will provide a separate option. Composition and color cues \citep{celona2022composition} will guide a few output candidates. A small model using frozen visual features \citep{oquab2023dinov2} will learn to select the candidate that best fits the requested look. This selection step supports the editing tool.

\section{Testing the Results}

I plan 120 new photos, split 80/20/20 for training/validation/testing by source. For each editing goal, I will compare fixed presets, subject-aware processing, and learned selection on shared candidates. Blinded author comparisons will assess whether outputs match the requested look, preserve subject details, and avoid visible artifacts. I will also measure processing time and examine failures such as poor masks. Online popularity guides reference selection, not the quality score.

\section{Current Progress and Schedule}

The prototype supports subject segmentation, composition controls, and foreground color preservation. Eight of 12 development images passed segmentation/geometric checks, producing 64 replay-verified candidates; this does not establish visual quality. Reference curation and data preparation are planned for October, comparisons and model training for November, and final reporting for early December.

{\small
\begin{thebibliography}{3}
\setlength{\bibsep}{0pt}
\bibitem{wang2024instantstyle}
H. Wang et al.
\href{https://arxiv.org/abs/2404.02733}{InstantStyle: Free Lunch towards Style-Preserving in Text-to-Image Generation}.
\emph{arXiv:2404.02733}, 2024.
\bibitem{celona2022composition}
L. Celona et al.
\href{https://arxiv.org/abs/2111.04647}{Composition and Style Attributes Guided Image Aesthetic Assessment}.
\emph{IEEE Transactions on Image Processing}, 31:5009-5024, 2022.
\bibitem{oquab2023dinov2}
M. Oquab et al.
\href{https://arxiv.org/abs/2304.07193}{DINOv2: Learning Robust Visual Features without Supervision}.
\emph{arXiv:2304.07193}, 2023.
\end{thebibliography}
}
\end{document}
